\begin{theorem}\label{th:another-closability}
    Пусть $X, Y$ -- банаховы и $A \in \LL(X, Y)$ такой, что $\im A^*$ -- сильно замкнут в $X^*$.
    Тогда $\im A$ -- замкнут в $Y$ и
    \[\im A = {}^\perp (\ker A^*), \quad \im A^* = \ker A^\perp.\]
\end{theorem}
\begin{proof}
    Пусть $Z = [\im A] \subset Y$.
    Тогда $Z$ -- замкнутое подпространство банахова пространства, а значит банахово в себе.
    Рассмотрим $T\colon X \to Z, x \mapsto Ax$.
    Так как $\im T = \im A$, то $\im T$ -- всюду плотен в $Z$.
    Рассмотрим $T^*\colon Z^* \to X^*$.
    По теореме Фредгольма 
    \[
        \ker T^* = (\im T)^\perp.
    \]
    Но $\im T^\perp = \{g \in Z^* \mid g(Tx) = 0, \forall x \in X\}$, а так как $\im T$ -- всюду плотен в $Z$ и $T$ -- непрерывен, то он состоит из функционалов, зануляющихся на всём $Z$ -- то есть только из нуля (в силу непрерывности и равенства нулю на плотном подмножестве).
    Поэтому ядро $T^*$ -- тривиально, а значит существует $(T^*)^{-1}\colon \im T^* \to Z^*$.

    Покажем, что $\im T^* = \im A^*$.
    Пусть $h \in \im T^*$. 
    Тогда существует $g \in Z^*$ такой, что $h = T^* g$.
    По определению 
    \[
        \forall x \in X \quad h(x) = (T^*g)(x) = g(Tx) = g(Ax).
    \]
    Продолжим $g\colon Z \to \Cm$ на всё $Y$ как $f\in Y^*$ по теореме Хана-Банаха с сохранением нормы. Тогда
    \[
        \forall x \in X \quad h(x) = g(Ax) = f(Ax) = (A^* f)(x).
    \]
    Поэтому $h = A^* f$ и, следовательно, $\im T^* \subset \im A^*$.
    Если же $h \in \im A^*$, то $\exists f \in Y^*\colon h = A^* f$.
    То есть
    \[
        h(x) = (A^*f)(x) = f(Ax) = f(Tx) = f|_{Z}(Tx) = (T^* f|_{Z})(x).
    \]
    Это показывает обратное вложение образов.

    Следовательно $\im T^* = \im A^*$ -- сильно замкнуто и, таким образом, полно (как замкнутое подпространство банахова $X^*$).
    По теореме Банаха об обратном операторе $\|(T^*)^{-1}\| < +\infty$.

    Заметим, что для $Z = \im A = \im T$ достаточно показать, что $\exists r > 0$ такой, что $O^Z_r(0) \subset A(O_1^X(0))$.
    Рассмотрим $M = [T(O_1^X(0))] \subset Z$.
    Это множество является выпуклым сильно замкнутым подмножеством.
    Для $z \in Z \setminus M$ воспользуется теоремой отделимости и получим, что $\exists g \in Z^*, \gamma \in \R$ такие, что 
    \[
        \forall x \in O_1^X(0)\hookleftarrow \Re g(z) > \gamma \geq \Re g(Tx)
    \]
    Так как 
    \[
        \Re g(Tx) = \Re (T^* g)(x),
    \]
    то если мы возьмём супремум по $x \in O_1^X(0)$, получится
    \[
        \Re g(z) > \gamma \geq \| \Re T^* g\| = \|T^* g\|.
    \]
    Но $g = (T^*)^{-1}T^* g$.
    То есть
    \[
        \Re (T^*)^{-1}(T^* g)(z) \geq \gamma \geq \|T^* g\|.
    \]
    Поскольку 
    \[
        \Re (T^*)^{-1}(T^* g)(z) \leq \|(T^*)^{-1}T^* g\|\|z\| \leq \|(T^*)^{-1}\|\|T^* g\|\|z\|,
    \]
    то 
    \[
        \|(T^*)^{-1}\|\|z\|\|T^* g\| > \gamma \geq \|T^*g\|.
    \]
    В частности, так как неравенство -- строгое, то $T^* g \neq 0$, и
    \[
        \|(T^*)^{-1}\|\|z\| \geq 1 \Rightarrow \|z\| \geq \frac{1}{\|(T^*)^{-1}\|} = r.
    \]
    Поэтому $O_r^Z(0) \subset B^Z_r(0) \subset [TO_1^X(0)]$.
    
    В частности, $\forall \epsilon > 0 \hookrightarrow O_{r\epsilon}^Z(0) \subset [TO_{\epsilon}^X(0)]$.
    Покажем, что $O_r^Z(0) \subset TO_3^X(0)$.

    Пусть $z \in [TO_1^X(0)]$.
    Так как $z - [TO_{1/2}^X(0)] \supset O_{r/2}^Z(z)$, то $z - [TO_{1/2}^X(0)] \cap TO_1^X(0) \neq \emptyset$. Поэтому существуют $z_1 \in [TO_{1/2}^X(0)]$ и $Tx_1 \in TO_{1}^X(0)$ такие, что $z - z_1 = y_1$.
    Аналогично, $z_1 - [TO_{1/4}^X(0)] \supset O_{r/4}^{Z}(z_1)$, а значит $z_1 - [TO_{1/4}^X(0)] \cap TO_{1/2}^X(0) \neq \emptyset$.
    И тогда существуют $z_2 \in [TO_{1/4}^X(0)]$ и $Tx_2 \in TO_{1/2}^X(0)$ такие, что $z_1 - z_2 = Tx_2$.
    Продолжая по индукции, получим $z_m \in [TO_{1/2^m}^X(0)]$ и $Tx_m \in TO^X_{1/2^m}(0)$ такие, что $z_m - z_{m + 1} = Tx_{m + 1}$.

    Но
    \[
        z - z_{m + 1} = (z - z_1) + \ldots + (z_m - z_{m + 1}) = T(x_1 + \ldots x_{m + 1}),
    \]
    а так как
    \[
        \sum_{k = 1}^{\infty} \|x_k\| \leq \sum_{k = 0}^{\infty} 2^{-k}  = 2,
    \]
    поэтому $\sum_{k = 1}^{m} x_k \rightarrow w \in X$, причём $\|w\| \leq 2$.
    А так как $z - z_{m + 1} \rightarrow z$, то $z = Tw$ -- что и даёт нам искомое утверждение.

    Для всякого $z \in Z \setminus \{0\}$ рассмотрим $\frac{z}{\|z\|}\frac{r}{2} \in O^Z_{r/3}(0) \subset A(O_1^X(0)) = T(O_1^X(0))$, то есть $\exists x \in O_1^X(0)\colon \frac{z}{\|z\|}\frac{r}{2} = Ax = Tx$ -- что и даёт $Z = \im A = \im T$.

    Так как $\im A$ -- замкнут, то из $\ker A^* = \im A^\perp$ получаем навершиванием аннулятора слева
    \[
        \im A = {}^\perp \ker A^*.
    \]
    Второе равенство получено ранее.
\end{proof}
\begin{example}
    Банаховость $X$ в условиях предыдущей теоремы существенна.
    Пусть $A\colon (c_{00}, \|\cdot\|_{\infty}) \to (c_0, \|\cdot\|_\infty), x \mapsto x$.
    Тогда $A^*\colon (\ell^1, \|\cdot\|_1) \to (\ell^1, \|\cdot\|_1)$, а также $A^*f = f$, и $\im A^* = \ell^1$ -- сильно замкнут в $\ell^1$.
    Но $\im A = c_{00}$ -- не является замкнутым в $c_0$ (а замыканием $c_{00}$ будет всё $c_0$).
\end{example}
\section{Компактные операторы.}
\begin{definition}
    Пусть $X, Y$ -- нормированные пространства.
    Будем называть оператор $A\colon X \to Y$ компактным, если $AB_1(0)$ -- вполне ограничен в $Y$.
\end{definition}
\begin{theorem}\label{th:finite-dimension-kernel}
    Рассмотрим $X$ --банахово пространство, $A \in \LL(X)$ -- компактный оператор, $I$ -- единичный оператор и $\lambda \in \Cm \setminus \{0\}$.

    Пусть $A_\lambda = A - \lambda I$.
    Тогда
    \begin{enumerate}
        \item $\dim\ker A_\lambda <+\infty$
        \item $\im A_\lambda$ -- замкнут в $X$.
    \end{enumerate}
\end{theorem}
\begin{proof}
    Покажем конечномерность ядра.
    Для этого достаточно показать, что единичный шар является вполне ограниченным, то есть для произвольной ограниченной последовательности $\{x_n\} \subset \ker A_\lambda\colon \|x_n\| \leq R$ существует фундаментальная подпоследовательность.

    Заметим, что $A_\lambda x_n = 0 \Longleftrightarrow x_n = \frac{1}{\lambda}Ax_n \in \frac{1}{\lambda} AB_R(0)$.
    Но в силу компактности $A$ множество $\frac{1}{\lambda}AB_R(0)$ -- вполне ограничено в $X$.
    И, следовательно, в последовательности $\{x_n\}$ существует фундаментальная подпоследовательность (что и требовалось).

    Теперь покажем замкнутость образа.
    Пусть $\dim \ker A_\lambda = n$ и $e_1, \ldots, e_n$ -- базис в $\ker A_\lambda$.
    То есть 
    \[
        \ker A_\lambda \ni x = \sum_{k = 1}^{n}\alpha_k(x)e_k,
    \]
    где $\alpha_k\colon \ker A_\lambda \to \Cm$ -- координатные функционалы.
    Они непрерывны (в силу конечномерности домена), а значит, согласно теореме Хана-Банаха, мы можем их продолжить на всё $X$ как $f_1, \ldots, f_n$ с сохранением нормы.
    Тогда $\forall x \in X$ верно, что $\sum_{k = 1}^n f_k(x)e_k = y \in \ker A_\lambda$,
    и, таким образом
    \[
        x - y \in \bigcap_{k = 1}^{n} \ker f_k.
    \]
    То есть мы получаем следующее разложение $X$:
    \[
        X = \ker A_\lambda \oplus \bigcap_{k = 1}^{n} \ker f_k = \ker A_\lambda \oplus N.
    \]
    
    Покажем, что $\exists R > 0$ такое, что $\forall x \in N \hookrightarrow \|A_\lambda(x)\| \geq R\|x\|$. Пусть это не так.
    Тогда $\forall R > 0 \exists x_R \in N\colon \|A_\lambda x_R\| < R\|x_R\|$.
    Возьмём $R_m = \frac{1}{m}$ и $z_m = \frac{x_{R_m}}{\|x_{R_m}\|} \in N$ (что эквивалентно тривиальности ядра $A$ на $N$).
    Норма образа оценивается как $\|A_\lambda z_m\| < \frac{1}{m}$.
    Так как $A$ -- компактный, то образ вполне ограничен и мы можем выделить фундаментальную подпоследовательность $\{Az_{m_k}\}$, которая в силу полноты $X$ сходится к некоторому $w \in X$.
    Но 
    \[
        z_{m_k} = \frac{A_\lambda z_{m_k} - Az_{m_k}}{-\lambda} \rightarrow \frac{w}{\lambda}.
    \]
    А значит и $z_{m_k}$ -- сходящаяся.
    Так как $\|A_\lambda \frac{w}{\lambda}\| = 0$, то $w \in \ker A_\lambda$
    При этом в силу замкнутости $N$ вектор $w \in N$.
    То есть $w = 0$. 
    Противоречие, так как норма $\|w\| = |\lambda|$.

    Тогда для всякого $y \in [\im A_\lambda]$ существует $x_m$ такая, что
    \[
        y = \lim_{m \rightarrow \infty} A_\lambda x_m.
    \]
    Представим $z_m = x_m - \sum_{k = 1}^{n} f_k(x_m)e_k$.
    То есть
    \[
        y = \lim A_\lambda z_m.
    \]
    А значит в силу наличия непрерывной обратимости на $N$
    \[
        \|z_m - z_n\| \leq \frac{1}{R}\|A_\lambda z_m - A_\lambda z_n\| \rightarrow 0,
    \]
    и тогда $z_m$ фундаментальна и, следовательно, сходящаяся последовательность -- то есть $y = A_\lambda z$.
\end{proof}
\begin{note}
    В условиях предыдущей теоремы $\im A_\lambda = ^\perp \ker (A_\lambda )^* = ^\perp\ker A^*_\lambda$, и $\im A_\lambda^* = (\ker A_\lambda)^\perp$.
    Поэтому $A_\lambda u = v$ (уравнение Фредгольма) -- разрешимо тогда и только тогда, когда $f(v) = 0$ при всяком $f \in \ker A_\lambda^*$.
    Аналогично и для <<союзного>> уравнения $A^*_\lambda f = g$.
\end{note}